\section{Level-Set Preservation in Output Space}

Section~\ref{sec:scalar-functional-preservation} preserves a scalar functional value directly. A second construction instead preserves the model output up to motions along a level set of a differentiable output functional $\phi_a:\mathcal Z\to\mathbb R$. This output-space perspective connects to distillation and functional regularization approaches to continual learning \citep{LwF,Distillation,FunctionalGP,FROMP}. Let $a \in D_A$ be an old sample, write $z_A=f_{\theta_A^\star}(x)$, and define the exact preservation set
\[
\mathcal C_a=
\{z:\phi_a(z)=\phi_a(z_A)\}.
\]
Thus the preserved quantity is the model output itself:
\[
q(a,\theta)=f_\theta(x).  
\]
The hard constraint is
\[
q(a,\theta)\in\mathcal C_a,
\]
and the corresponding unscaled proximal loss is
\[
p(a,\theta) =
\operatorname{dist}_{G}\big(q(a,\theta),\mathcal C_a\big)^2,
\]
where $G\succ0$ is the chosen local output metric. Positive definiteness makes the primal metric and its inverse on covectors unambiguous; Section~\ref{sec:output-geometry} treats the restricted pseudoinverse convention needed by a singular ordered-output example.

We obtain a tractable penalty from the standard first-order model of the constraint set near $z_A$, namely its tangent plane. Let
\[
r=
\nabla_z \phi_a(z)\big|_{z=z_A}.
\]
We first assume $r\ne0$, so the functional's level set has a regular tangent hyperplane at $z_A$.
For a small output displacement $\Delta z=f_\theta(x)-f_{\theta_A^\star}(x)$, its first-order change is approximately $r^\top \Delta z$.
Hence the local constraint set is
\[
T=
\{\Delta z:r^\top \Delta z=0\}.  
\]
Replacing the exact level set by this tangent hyperplane gives the local proximal loss
\[
p(a,\theta)  
\approx  
\operatorname{dist}_{G}(\Delta z,T)^2.  
\]

For the metric $|v|_G^2=v^\top G v$, the squared distance to the tangent hyperplane follows from projecting $\Delta z$ onto the gradient:
\[
\operatorname{dist}_{G}(\Delta z,T)^2 =
\frac{(r^\top \Delta z)^2}{r^\top G^{-1}r}.  
\]
To verify the matched metric pair, write the normal residual as $e=\Delta z-t$ with $t\in T$. It solves
\[
\min_e e^\top G e
\quad\text{subject to}\quad
r^\top e=r^\top\Delta z.
\]
The Lagrange equations give
\[
e^\star
=
\frac{r^\top\Delta z}{r^\top G^{-1}r}G^{-1}r,
\]
whose squared $G$-norm is the displayed quotient. Thus $G$ measures primal output displacement while $G^{-1}$ acts on the loss covector; these roles cannot be chosen independently.
Linearizing the model output around $\theta_A^\star$, with
\[
\delta=\theta-\theta_A^\star,  
\qquad  
J=\nabla_\theta f_{\theta_A^\star}(x),  
\]
we have
\[
\Delta z \approx J\delta.  
\]
Substituting this into the local regularization term gives
\[
\frac{\lambda}{2}p(a,\theta_A^\star+\delta)  
\approx  
\frac{\lambda}{2}  
\frac{(r^\top J\delta)^2}{r^\top G^{-1}r}.  
\]
Equivalently,
\[
\frac{\lambda}{2}p(a,\theta_A^\star+\delta)  
\approx  
\frac{\lambda}{2}  
\delta^\top  
\left(
J^\top  
\frac{rr^\top}{r^\top G^{-1}r}  
J  
\right)
\delta.  
\]

Taking the empirical average over old-distribution samples $a_1,\dots,a_N$ yields an importance matrix of the form
\[
M_{2,A} =
\frac{1}{N}\sum_{i=1}^N  
J_i^\top  
\frac{r_i r_i^\top}{r_i^\top G_i^{-1}r_i}  
J_i.  
\]

This local construction penalizes motion normal to the tangent model. IEWC is the evaluated specialization $\phi_a=\ell_A$; other functionals define CLOC extensions. When $r_i=0$, the tangent normal is undefined and the undamped quotient is $0/0$; our finite-damping implementation assigns that sample zero contribution. Section~\ref{sec:well-fitted} characterizes the resulting near-interpolation behavior.
